\input{../tex-inputs/latex-leseansicht-vorspann}

\section[Carl von Torresani an Arthur Schnitzler, {[}22. 4. 1906{]}]{L04324 Carl von Torresani an Arthur Schnitzler, {[}22. 4. 1906{]}}
\nopagebreak\mylabel{L04324v}
\rehead{ }\normalsize\beginnumbering\briefempfaengerindex{Schnitzler, Arthur@\textsc{Schnitzler, Arthur}!zzzTorresani-Lanzenfeld, Carl von@\emph{von Carl von Torresani-Lanzenfeld}!1906-04-222@{{[}22. 4. 1906{]}}|(be}
\toendnotes[C]{\smallbreak\pagebreak[2]}
\correspDesc{Versand  durch Carl von Torresani am [22. 4. 1906] in Torbole sul Garda
\newline{}Erhalt  durch Arthur Schnitzler am [23. 4. 1906] in Wien}\toendnotes[C]{\smallbreak}
\Standort{CUL, Schnitzler, B 106.}
\physDesc{Telegramm, 125 Zeichen
\newline{}maschinell
\newline{}Schnitzler: mit Bleistift datiert: »22/4 90\textcolor{gray}{6}« }\toendnotes[C]{\smallbreak}
\pstart
           \centering{}{\pb}fr torbole\oindex{Torbole sul Garda@\textbf{Torbole sul Garda}, \emph{Hauptstadt}|pw} 223 19 23{ }12/10 n\pend
           \vspace{0.5em}
\pstart
           hocherfreut und geruehrt durch ihr gedenken \label{K_L04324-1v}\edtext{sende innigen Dank}{\lemma{\textnormal{\emph{sende innigen Dank}}}\Cendnote{\textnormal{Schnitzlers Gruß ist nicht überliefert. Es dürfte sich es sich um eine Gratulation zu Torresanis\pwindex{Torresani-Lanzenfeld, Carl von 19.\,4.\,1846 Mailand – 16.\,4.\,1907 Nago-Torbole@\textsc{Torresani-Lanzenfeld, Carl von} (19.\,4.\,1846 Mailand – 16.\,4.\,1907 Nago-Torbole), \emph{Schriftsteller, Offizier}|pwk} 60. Geburtstag am 19. 4. 1906 gehandelt haben, der mit militärischen Feierlichkeiten in Torbole\oindex{Torbole sul Garda@\textbf{Torbole sul Garda}, \emph{Hauptstadt}|pwk} am Gardasee\oindex{Lago di Garda@\textbf{Lago di Garda}, \emph{See}|pwk} begangen wurde. Über die vielen Glückwunschadressen, unter deren Absendern auch »viele Schriftsteller« waren, berichtete sogar die Zeitung, siehe \emph{(Karl Baron von Torresani)}. In: \emph{Wiener Zeitung}\pwindex{Wiener Zeitung@\emph{Wiener Zeitung}|pwk}, Nr. 97, 28. 4. 1906, S. 8.}}}\label{K_L04324-1} in bewunderung und
               ergebenst =\pend
           \pstart \spacefill\mbox{torresani}\pend{}\selectlanguage{ngerman}\endnumbering\briefempfaengerindex{Schnitzler, Arthur@\textsc{Schnitzler, Arthur}!zzzTorresani-Lanzenfeld, Carl von@\emph{von Carl von Torresani-Lanzenfeld}!1906-04-222@{{[}22. 4. 1906{]}}|)be}\mylabel{L04324h}
\begin{anhang}
\end{anhang}\newcommand{\dateiname}{L04324}\newcommand{\titel}{Carl von Torresani an Arthur Schnitzler, [22. 4. 1906]}\newcommand{\editorInnen}{Herausgegeben von Jahnke, SelmaMüller, Martin Anton}\input{../tex-inputs/latex-leseansicht-abspann}
